\chapter{Introduction\label{cha:introduction}}
% Drafted per POINT-SEQUENCE.md Ch.1 (points 1.1–1.12).
% Note: Intro precedes Background, so every technical term here is used in its plain-word form
% and gets its precise definition in Chapter 2.

\section{Motivation\label{sec:motivation}}

% [1.1]
Sensors in cities, factories, and vehicles produce data without pause.
Stream processing systems compute results from this data while it arrives.

% [1.2]
Many useful computations need memory of the recent past.
A system that reports the average temperature per sensor and per minute must collect values until each minute is complete.
These collected partial results are called state.

% [1.3]
NebulaStream is a stream processing system for the Internet of Things~\cite{zeuch2020,zeuch2022}.
It is built to run on small computers near the sensors, at the edge of the network.
NebulaStream keeps all state in main memory.

% [1.4]
Edge devices have small main memories.
A typical device offers a few hundred megabytes to a few gigabytes.

% [1.5]
State grows with the number of sensors and the length of the observed time range.
When state grows beyond the available memory, the operating system terminates the process.
Processing halts and results are lost.

% [1.6]
Two common answers to this problem fail at the edge.
Provisioning more memory is not possible on fixed hardware.
Moving the computation to the cloud forces the raw data to travel, and data movement dominates the cost of edge deployments~\cite{horowitz2014,zeuch2020}.

% [1.7]
Cloud stream processing systems know a third answer.
An out-of-core state backend moves state to local disk when memory runs short.
Out-of-core means that a computation uses more data than fits into main memory.
RocksDB, an embeddable storage engine, is the industry standard for this role~\cite{flink2021}.

\section{Research Challenge\label{sec:challenge}}

% [1.8]
NebulaStream executes queries with a push-based model.
Data is pushed through the system as it arrives, and all operators share one fixed pool of memory buffers.
Whether an out-of-core backend works inside this model is unknown.
The combination has not been evaluated systematically~\cite{marcu2022}.

% [1.9]
This thesis builds an out-of-core backend for window state in NebulaStream.
It then evaluates the backend on edge-class memory budgets, from 512 megabytes to 2 gigabytes.

\section{Results in Brief\label{sec:results-brief}}

% [1.10]
The backend processes workloads whose state exceeds memory, without crashes, inside an operating window that this thesis characterizes.
The throughput cost ranges from 3\% to 22\%.
During the evaluation, one engine configuration default turned out to silently disable the engine's intended memory architecture.
Correcting it requires one flag.
The correction, combined with new memory-pressure mechanisms, raises stability to the level of the in-memory engine at 25\% less resident memory.
The one remaining failure mode is diagnosed precisely, and its fix is specified as future work.

\section{Contributions\label{sec:contributions}}

% [1.11]
This thesis makes six contributions.
\begin{enumerate}
  \item \textbf{C1 — A spillable window state backend for NebulaStream.}
    Window state moves to RocksDB under memory pressure and returns for result emission.
    The query compiler and the just-in-time compilation pipeline remain untouched.
  \item \textbf{C2 — An instrumented, reproducible evaluation harness.}
    Every run executes under a hard memory cap with pinned processor cores.
    Six probes record the memory system over time.
    The evaluation covers roughly 3{,}100 measured runs, and every follow-up hypothesis was registered before its experiment ran.
  \item \textbf{C3 — Empirical operating windows for out-of-core processing on edge budgets.}
    The evaluation maps where the backend succeeds and where it fails, per memory budget.
    It shows that completion is a race between data arrival and disk drainage, not a fixed state-size threshold.
  \item \textbf{C4 — A load-bearing architecture finding.}
    The engine's default buffer size silently kept all window state outside the engine's buffer pool in every configuration.
    One flag corrects this.
    The corrected architecture, plus the new pressure mechanisms, reaches in-memory stability at lower memory cost.
  \item \textbf{C5 — A failure taxonomy with a diagnosed root cause.}
    Failures follow a three-regime law.
    The residual failure is an allocation deadlock, demonstrated by direct instrumentation.
  \item \textbf{C6 — An open-source implementation.}
    All mechanisms are configuration-gated.
    With all gates closed, the engine is byte-identical to the unmodified baseline.
\end{enumerate}

\section{Thesis Structure\label{sec:structure}}

% [1.12]
Chapter~\ref{cha:background} introduces NebulaStream, windowed aggregation, RocksDB, and the terms of the evaluation.
Chapter~\ref{cha:problem} states the research problem and the research questions.
Chapter~\ref{cha:solution} presents the design and implementation of the backend.
Chapter~\ref{cha:evaluation} evaluates effectiveness, cost, and stability, and investigates the failure modes.
Chapter~\ref{cha:related} relates the thesis to existing work.
Chapter~\ref{cha:conclusion} concludes and lists future work.
